% ============================================================
% APUNTES.pdf -- Trozo 1
% Páginas 1--2
% ============================================================

\section{Programación lineal y optimización convexa}

\subsection{Contenido teórico del bloque}

En esta parte del curso aparecen, en términos generales, los siguientes temas:

\begin{itemize}
    \item \textbf{Programas lineales.}
    \begin{itemize}
        \item Conceptos básicos y ejemplos.
        \item Teoría de programación lineal.
        \item Método simplex.
        \item Dualidad y lema de Farkas.
    \end{itemize}

    \item \textbf{Programas cónicos.}

    \item \textbf{Programas semidefinidos.}

    \item \textbf{Optimización convexa.}
    \begin{itemize}
        \item Conceptos básicos.
        \item Teoremas de optimalidad de primer y segundo orden.
        \item Condiciones KKT.
        \item Dualidad.
    \end{itemize}
\end{itemize}

\subsection{Problemas de optimización}

Un problema de optimización tiene, en general, la forma
\[
    \min_{x} f(x)
    \qquad \text{sujeto a } x\in \Omega.
\]
También se escribe de manera compacta como
\[
    \min\{f(x):x\in \Omega\}.
\]

Aquí:
\begin{itemize}
    \item \(x\) es la \emph{variable de decisión};
    \item \(f\) es la \emph{función objetivo};
    \item \(\Omega\) es el \emph{conjunto factible}.
\end{itemize}

Los componentes de \(x\), es decir,
\[
    x=(x_1,\dots,x_n),
\]
se denominan \emph{variables de optimización} o \emph{variables de decisión}.

\begin{definition}[Solución óptima]
Sea \(\Omega\subseteq \mathbb{R}^n\). Un punto \(x^\star\in \Omega\) se llama \emph{solución óptima} o \emph{minimizador} del problema
\[
    \min\{f(x):x\in \Omega\}
\]
si
\[
    f(x^\star)\leq f(x)
    \qquad \forall x\in \Omega.
\]
\end{definition}

En otras palabras, \(x^\star\) es factible y alcanza el menor valor posible de la función objetivo sobre el conjunto factible.

\subsection{Restricciones funcionales}

En muchos problemas de optimización, el conjunto factible se describe mediante restricciones de desigualdad e igualdad. Una forma típica es
\[
    \Omega
    =
    \left\{
    x\in F:
    \begin{array}{ll}
        g_j(x)\leq 0, & j=1,\dots,r,\\[2mm]
        h_i(x)=0, & i=1,\dots,m.
    \end{array}
    \right\}.
\]

En este caso:
\begin{itemize}
    \item las funciones \(g_j\) representan restricciones de desigualdad;
    \item las funciones \(h_i\) representan restricciones de igualdad;
    \item \(F\) representa un conjunto ambiente, usualmente \(F\subseteq \mathbb{R}^n\).
\end{itemize}

\subsection{Optimización convexa y programación lineal}

\subsubsection{Optimización convexa}

Un problema de optimización convexa tiene la forma
\[
    \min f(x)
    \qquad
    \text{sujeto a } x\in \Omega,
\]
donde la estructura convexa aparece tanto en la función objetivo como en las restricciones.

De manera más explícita, en un problema convexo con restricciones funcionales, se pide que:
\begin{itemize}
    \item la función objetivo \(f\) sea convexa;
    \item las funciones de desigualdad \(g_j\) sean convexas;
    \item las funciones de igualdad \(h_i\) sean afines.
\end{itemize}

Es decir, una forma estándar es
\[
    \begin{array}{ll}
        \min & f(x)\\[1mm]
        \text{sujeto a} & g_j(x)\leq 0,\qquad j=1,\dots,r,\\
        & h_i(x)=0,\qquad i=1,\dots,m,
    \end{array}
\]
con \(f\) y \(g_j\) convexas, y \(h_i\) afines.

\subsubsection{Programación lineal}

Un \emph{programa lineal} corresponde al caso en que la función objetivo y las restricciones son afines en \(x\). Es decir,
\[
    f,\ g_j,\ h_i
    \quad
    \text{son funciones afines en } x.
\]

Recordemos que una función afín tiene la forma
\[
    a^\top x+b,
\]
donde \(a\in\mathbb{R}^n\) y \(b\in\mathbb{R}\).

Por tanto, un programa lineal puede escribirse en una forma del tipo
\[
    \begin{array}{ll}
        \min & c^\top x\\[1mm]
        \text{sujeto a} & Ax\leq b,\\
        & Ex=d,
    \end{array}
\]
donde las restricciones son lineales o afines.

\subsection{Existencia de soluciones}

\subsubsection{Teorema de Weierstrass}

\begin{theorem}[Weierstrass]
Sea \(S\subseteq\mathbb{R}^n\) un conjunto no vacío y compacto. Si
\[
    f:S\to\mathbb{R}
\]
es continua, entonces \(f\) alcanza su mínimo y su máximo sobre \(S\).
\end{theorem}

En particular, el problema
\[
    \min\{f(x):x\in S\}
\]
tiene solución siempre que \(S\) sea no vacío y compacto, y \(f\) sea continua.

\subsubsection{Funciones sin minimizador}

La continuidad de \(f\) no basta por sí sola para garantizar existencia de minimizador. El conjunto factible también importa.

\paragraph{Ejemplo 1.}
Considérese
\[
    f(x)=\tan(x),
    \qquad
    \Omega=\left(-\frac{\pi}{2},\frac{\pi}{2}\right).
\]
Entonces \(f\) es continua en \(\Omega\), pero
\[
    \tan(x)\to -\infty
    \qquad
    \text{cuando } x\to -\frac{\pi}{2}^{+}.
\]
Por lo tanto, los valores de \(f\) pueden ser arbitrariamente pequeños y el problema
\[
    \min\{\tan(x):x\in \Omega\}
\]
no tiene minimizador.

\paragraph{Ejemplo 2.}
Considérese
\[
    f(x)=C+e^x,
    \qquad
    \Omega=\mathbb{R}.
\]
Entonces
\[
    C+e^x>C
    \qquad
    \forall x\in\mathbb{R},
\]
pero
\[
    C+e^x\to C
    \qquad
    \text{cuando } x\to -\infty.
\]
Así,
\[
    \inf_{x\in\mathbb{R}} f(x)=C,
\]
pero no existe \(x\in\mathbb{R}\) tal que \(f(x)=C\). Por tanto, el mínimo no se alcanza.

\subsection{Subniveles y existencia de minimizadores}

\begin{definition}[Conjunto de subnivel]
Sea \(f:S\to\mathbb{R}\). Para \(\alpha\in\mathbb{R}\), el conjunto de subnivel asociado a \(\alpha\) se define por
\[
    L_\alpha(f)
    =
    \{x\in S:f(x)\leq \alpha\}.
\]
\end{definition}

\begin{theorem}[Criterio de existencia mediante subniveles]
Sea \(f:S\to\mathbb{R}\) continua. Si existe \(\alpha\in\mathbb{R}\) tal que
\[
    L_\alpha(f)=\{x\in S:f(x)\leq \alpha\}
\]
es no vacío y compacto, entonces \(f\) tiene un mínimo en \(S\).
\end{theorem}

\begin{proof}
Como \(L_\alpha(f)\) es compacto y \(f\) es continua, por el teorema de Weierstrass existe
\[
    x^\star\in L_\alpha(f)
\]
tal que
\[
    f(x^\star)\leq f(x)
    \qquad
    \forall x\in L_\alpha(f).
\]
Además, si \(x\notin L_\alpha(f)\), entonces
\[
    f(x)>\alpha.
\]
Como \(x^\star\in L_\alpha(f)\), se tiene \(f(x^\star)\leq \alpha\). Por tanto,
\[
    f(x^\star)\leq \alpha < f(x)
    \qquad
    \forall x\notin L_\alpha(f).
\]
Luego,
\[
    f(x^\star)\leq f(x)
    \qquad
    \forall x\in S.
\]
Así, \(x^\star\) es minimizador global de \(f\) sobre \(S\).
\end{proof}

\subsection{Ejemplo: problema de proyección}

Sea \(y\in\mathbb{R}^n\) y sea \(S\subseteq\mathbb{R}^n\) un conjunto cerrado, no necesariamente acotado. Queremos encontrar un punto de \(S\) que esté lo más cerca posible de \(y\). Esto conduce al problema
\[
    \min_{x\in S}\|x-y\|.
\]

Equivalentemente,
\[
    \begin{array}{ll}
        \min & f(x)=\|x-y\|\\[1mm]
        \text{sujeto a} & x\in S.
    \end{array}
\]

Aunque \(S\) no sea compacto, el problema tiene solución. La idea es restringir el problema a un conjunto compacto que contiene algún minimizador del problema original.

Como \(S\neq\varnothing\), tomamos un punto \(z\in S\). Entonces basta considerar los puntos \(x\in S\) que satisfacen
\[
    \|x-y\|\leq \|z-y\|.
\]
En efecto, cualquier punto \(x\in S\) con
\[
    \|x-y\|>\|z-y\|
\]
no puede ser minimizador, porque \(z\) ya entrega un valor menor.

Por tanto, podemos restringir el problema al conjunto
\[
    \Omega
    =
    S\cap \overline{B}\bigl(y,\|z-y\|\bigr).
\]
Este conjunto es no vacío, cerrado y acotado; luego es compacto. Además, la función
\[
    f(x)=\|x-y\|
\]
es continua. Por Weierstrass, existe \(x^\star\in\Omega\) tal que
\[
    \|x^\star-y\|
    =
    \min_{x\in \Omega}\|x-y\|.
\]
Por la construcción de \(\Omega\), este mismo punto es minimizador del problema original sobre \(S\).

\subsection{El simplex de probabilidades}

Un conjunto importante en optimización es el simplex
\[
    \Delta_n
    =
    \left\{
    x=(x_1,\dots,x_n)\in\mathbb{R}_{\geq 0}^n:
    \sum_{i=1}^{n}x_i=1
    \right\}.
\]

Este conjunto puede interpretarse como el conjunto de distribuciones de probabilidad sobre el universo finito
\[
    \{1,\dots,n\}.
\]

En efecto, si \(x\in\Delta_n\), entonces:
\[
    x_i\geq 0
    \qquad
    \text{para todo } i=1,\dots,n,
\]
y
\[
    \sum_{i=1}^{n}x_i=1.
\]

\begin{remark}
El simplex \(\Delta_n\) es compacto. En efecto:
\begin{itemize}
    \item es cerrado, porque está definido mediante las condiciones cerradas
    \[
        x_i\geq 0,
        \qquad
        \sum_{i=1}^{n}x_i=1;
    \]
    \item es acotado, porque si \(x\in\Delta_n\), entonces
    \[
        0\leq x_i\leq 1
        \qquad
        \forall i=1,\dots,n.
    \]
\end{itemize}
Por tanto, \(\Delta_n\) es compacto.
\end{remark}

\subsection{Ejemplo: respuesta regularizada sobre el simplex}

Sea \(v=(v_1,\dots,v_n)\in\mathbb{R}^n\) una asignación de valores. Un problema típico consiste en elegir una distribución \(x\in\Delta_n\) que minimice una función compuesta por un término lineal y un término de regularización entrópica.

Un modelo de este tipo es
\[
    \min_{x\in\Delta_n}
    g(x)
    =
    -\sum_{i=1}^{n}v_i x_i
    +
    \frac{1}{\lambda}
    \sum_{i=1}^{n}x_i\log(x_i).
\]
% revisar: en el manuscrito aparece una constante de regularización;
% la forma más probable es (1/lambda) sum_i x_i log(x_i).

Equivalentemente,
\[
    \begin{array}{ll}
        \min & g(x)
        =
        -\displaystyle\sum_{i=1}^{n}v_i x_i
        +
        \dfrac{1}{\lambda}
        \displaystyle\sum_{i=1}^{n}x_i\log(x_i)
        \\[4mm]
        \text{sujeto a} & x\in \Delta_n.
    \end{array}
\]

Este tipo de problema aparece, por ejemplo, en versiones regularizadas de problemas de decisión o de mejor respuesta.

\subsection{Complejidad computacional}

Los problemas de programación lineal se resuelven eficientemente, tanto desde el punto de vista teórico como práctico.

\begin{remark}
En términos generales, la programación lineal pertenece a una clase de problemas tratables computacionalmente. En cambio, muchos problemas de optimización generales, especialmente cuando tienen restricciones discretas o combinatorias, pueden ser NP-duros.
\end{remark}

% revisar: en la página aparece una observación sobre P.L. y NP-duro.
% La interpretación más razonable es distinguir programación lineal continua
% de problemas generales de optimización o programación lineal entera.

% ============================================================
% APUNTES.pdf -- Trozo 2
% Páginas 3--4
% ============================================================

\subsection{Programación lineal entera y problema de corte máximo}

En general, los problemas de programación lineal entera pueden ser computacionalmente difíciles. Un ejemplo clásico es el problema de \emph{corte máximo}.

Sea \(G=(V,E)\) un grafo no dirigido. Dado un subconjunto \(S\subseteq V\), el corte asociado a \(S\) separa los vértices en dos partes:
\[
    S
    \qquad \text{y} \qquad
    V\setminus S.
\]

Una arista \(\{u,v\}\in E\) pertenece al corte si uno de sus extremos está en \(S\) y el otro está en \(V\setminus S\).

Para modelar esto, asignamos a cada vértice \(v\in V\) una variable
\[
    X_v\in\{-1,1\},
\]
con el siguiente significado:
\[
    X_v=1
    \quad \Longleftrightarrow \quad
    v\in S,
\]
y
\[
    X_v=-1
    \quad \Longleftrightarrow \quad
    v\in V\setminus S.
\]

Entonces, para una arista \(\{u,v\}\in E\), se tiene:
\[
    \{u,v\} \text{ cruza el corte}
    \quad \Longleftrightarrow \quad
    X_uX_v=-1.
\]
En cambio, si ambos extremos están en el mismo lado del corte, entonces
\[
    X_uX_v=1.
\]

Por tanto, el indicador de que la arista \(\{u,v\}\) cruza el corte puede escribirse como
\[
    \frac{1-X_uX_v}{2}.
\]

En consecuencia, el tamaño del corte está dado por
\[
    c(S,V\setminus S)
    =
    \sum_{\{u,v\}\in E}
    \frac{1-X_uX_v}{2}.
\]

Así, el problema de corte máximo puede formularse como
\[
    \begin{array}{ll}
        \max & \displaystyle
        \sum_{\{u,v\}\in E}
        \frac{1-X_uX_v}{2}
        \\[4mm]
        \text{sujeto a} & X_v\in\{-1,1\},
        \qquad v\in V.
    \end{array}
\]

De manera equivalente, usando la condición
\[
    X_v\in\{-1,1\}
    \quad \Longleftrightarrow \quad
    X_v^2=1,
\]
se puede escribir
\[
    \begin{array}{ll}
        \max & \displaystyle
        \sum_{\{u,v\}\in E}
        \frac{1-X_uX_v}{2}
        \\[4mm]
        \text{sujeto a} & X_v^2=1,
        \qquad v\in V.
    \end{array}
\]

\begin{remark}
Esta formulación ya no es un programa lineal continuo: aparecen restricciones cuadráticas y variables discretas. Esto refleja la dificultad combinatoria del problema.
\end{remark}

\subsection{Ejemplo geométrico de programa lineal}

Consideremos el problema lineal
\[
    \begin{array}{ll}
        \min & -x_1-x_2 \\[1mm]
        \text{sujeto a} & x_1\geq 0,\\
        & x_2\geq 0,\\
        & x_2-x_1\leq 1,\\
        & x_1+6x_2\leq 1,\\
        & 4x_1-x_2\leq 10.
    \end{array}
\]

Minimizar \(-x_1-x_2\) es equivalente a maximizar \(x_1+x_2\). Por tanto, el problema anterior puede leerse como
\[
    \max x_1+x_2
\]
sobre la región factible determinada por las desigualdades.

En forma matricial, el problema se puede escribir como
\[
    \begin{array}{ll}
        \min & f(x)=-x_1-x_2\\[1mm]
        \text{sujeto a}
        &
        \begin{pmatrix}
            -1 & 1\\
            1 & 6\\
            4 & -1
        \end{pmatrix}
        \begin{pmatrix}
            x_1\\
            x_2
        \end{pmatrix}
        \leq
        \begin{pmatrix}
            1\\
            1\\
            10
        \end{pmatrix},
        \\[6mm]
        & x_1\geq 0,\quad x_2\geq 0.
    \end{array}
\]

La pregunta geométrica es:

\[
    \text{¿Qué punto del polígono factible maximiza } x_1+x_2?
\]

La función objetivo
\[
    x_1+x_2
\]
tiene como conjuntos de nivel las rectas
\[
    N_\alpha
    =
    \{(x_1,x_2)\in\mathbb{R}^2:x_1+x_2=\alpha\}.
\]

Maximizar \(x_1+x_2\) equivale a desplazar estas rectas en la dirección del vector
\[
    (1,1)
\]
hasta encontrar el mayor nivel \(\alpha\) que todavía intersecta la región factible.

En otras palabras, una solución óptima está dada por el mayor valor de \(\alpha\) tal que
\[
    N_\alpha\cap \Omega\neq\varnothing,
\]
donde \(\Omega\) denota el conjunto factible.

\subsection{Formas equivalentes de problemas lineales}

En programación lineal, un mismo problema puede escribirse de distintas maneras equivalentes.

\subsubsection{Reemplazar igualdades por desigualdades}

Una igualdad puede reemplazarse por dos desigualdades. En efecto,
\[
    a^\top x=b
\]
es equivalente a imponer simultáneamente
\[
    a^\top x\leq b
    \qquad \text{y} \qquad
    -a^\top x\leq -b.
\]

Por tanto, una restricción de igualdad puede escribirse como dos restricciones de desigualdad.

\subsubsection{Invertir la dirección de optimización}

También podemos cambiar de maximización a minimización. Se tiene la equivalencia
\[
    \max c^\top x
    \quad \Longleftrightarrow \quad
    \min -c^\top x.
\]

De manera análoga,
\[
    \min c^\top x
    \quad \Longleftrightarrow \quad
    \max -c^\top x.
\]

\subsubsection{Forma de desigualdad}

Un programa lineal en forma de desigualdad puede escribirse como
\[
    \begin{array}{ll}
        \min & c^\top x\\[1mm]
        \text{sujeto a} & Ax\leq b.
    \end{array}
\]

Si al invertir direcciones o modificar restricciones se obtiene un conjunto factible vacío,
\[
    \Omega=\varnothing,
\]
entonces el problema no tiene solución factible.

\begin{definition}[Problema no factible]
Un problema de optimización se llama \emph{no factible} si su conjunto factible es vacío.
\end{definition}

\subsection{Ejemplo numérico de programa lineal}

En los apuntes aparece un ejemplo numérico con tres variables, escrito matricialmente como
\[
    \begin{array}{ll}
        \min & 7x_1+5x_2+15x_3\\[1mm]
        \text{sujeto a}
        &
        \begin{pmatrix}
            55 & 0.5 & 9.5\\
            60 & 300 & 10\\
            30 & 20 & 10
        \end{pmatrix}
        \begin{pmatrix}
            x_1\\
            x_2\\
            x_3
        \end{pmatrix}
        \leq
        \begin{pmatrix}
            0.5\\
            15\\
            4
        \end{pmatrix}.
    \end{array}
\]
% revisar: los coeficientes de esta matriz aparecen poco nítidos en la página 3.
% Especialmente el primer renglón y el término independiente 0.5 podrían requerir revisión.

\subsection{Ejemplo: problema de transporte}

Supongamos que hay dos fuentes o fábricas,
\[
    F_1,\quad F_2,
\]
y tres destinos,
\[
    D_1,\quad D_2,\quad D_3.
\]

Los costos unitarios de transporte están dados por la tabla
\[
\begin{array}{c|cc}
    & F_1 & F_2\\
    \hline
    D_1 & 2 & 1\\
    D_2 & 1 & 2\\
    D_3 & 3 & 1
\end{array}
\]

Definimos
\[
    x_{ij}
\]
como la cantidad enviada desde la fuente \(F_i\) hacia el destino \(D_j\), donde
\[
    i\in\{1,2\},
    \qquad
    j\in\{1,2,3\}.
\]

El costo total de transporte es
\[
    2x_{11}+x_{21}
    +
    x_{12}+2x_{22}
    +
    3x_{13}+x_{23}.
\]

Por tanto, el problema de transporte se puede escribir como
\[
    \begin{array}{ll}
        \min
        &
        2x_{11}+x_{21}
        +
        x_{12}+2x_{22}
        +
        3x_{13}+x_{23}
        \\[2mm]
        \text{sujeto a}
        &
        \displaystyle
        \sum_{j=1}^{3}x_{1j}=10000,
        \\[3mm]
        &
        \displaystyle
        \sum_{j=1}^{3}x_{2j}=5000,
        \\[3mm]
        &
        \displaystyle
        \sum_{i=1}^{2}x_{ij}=5000,
        \qquad j=1,2,3,
        \\[3mm]
        &
        x_{ij}\geq 0,
        \qquad i=1,2,\quad j=1,2,3.
    \end{array}
\]

Las dos primeras restricciones representan la capacidad o disponibilidad de cada fuente. Las tres restricciones siguientes representan la demanda de cada destino.

\subsection{Geometría de la programación lineal}

Una forma estándar importante de un programa lineal es la \emph{forma de igualdad}:
\[
    \begin{array}{ll}
        \max & c^\top x\\[1mm]
        \text{sujeto a} & Ax=b,\\
        & x\geq 0.
    \end{array}
\]

Aquí
\[
    A\in\mathbb{R}^{m\times n},
    \qquad
    m\leq n,
    \qquad
    \operatorname{rango}(A)=m.
\]

La pregunta central es:

\[
    \text{¿Cómo resolver } Ax=b \text{ con } x\geq 0?
\]

El conjunto factible asociado es
\[
    F
    =
    \{x\in\mathbb{R}^n:Ax=b,\ x\geq 0\}.
\]

\subsection{Variables de holgura}

Una desigualdad de tipo
\[
    a^\top x\leq b
\]
puede transformarse en una igualdad agregando una variable de holgura.

Por ejemplo, consideremos el problema
\[
    \begin{array}{ll}
        \max & 3x_1-2x_2\\[1mm]
        \text{sujeto a}
        &
        2x_1-x_2\leq 4,\\
        &
        -x_1+3x_2\leq -5,\\
        &
        x_2\geq 0.
    \end{array}
\]

Para convertir la restricción
\[
    2x_1-x_2\leq 4
\]
en una igualdad, añadimos una variable de holgura \(x_3\geq 0\), obteniendo
\[
    2x_1-x_2+x_3=4,
    \qquad
    x_3\geq 0.
\]

En general,
\[
    a^\top x\leq b
\]
se transforma en
\[
    a^\top x+s=b,
    \qquad
    s\geq 0.
\]

\subsection{Conjuntos convexos}

\begin{definition}[Conjunto convexo]
Un conjunto \(X\subseteq\mathbb{R}^n\) es \emph{convexo} si para cada par de puntos \(x,y\in X\), el segmento que une \(x\) con \(y\) también pertenece a \(X\). Es decir,
\[
    tx+(1-t)y\in X
    \qquad
    \forall x,y\in X,\quad \forall t\in[0,1].
\]
\end{definition}

\begin{proposition}
El conjunto de factibilidad de un programa lineal es convexo.
\end{proposition}

\begin{proof}
Tarea.
\end{proof}

\subsection{Conos}

\begin{definition}[Cono]
Un conjunto \(X\subseteq\mathbb{R}^n\) es un \emph{cono} si
\[
    \alpha x\in X
    \qquad
    \forall x\in X,\quad \forall \alpha\geq 0.
\]
\end{definition}

\begin{definition}[Cono convexo generado]
Sean
\[
    a_1,\dots,a_n\in\mathbb{R}^m.
\]
El \emph{cono convexo generado} por \(a_1,\dots,a_n\) es
\[
    \operatorname{cone}(a_1,\dots,a_n)
    =
    \left\{
        \sum_{i=1}^{n}t_i a_i:
        t_i\geq 0,\ i=1,\dots,n
    \right\}.
\]
\end{definition}

\begin{proposition}
El cono convexo generado por un conjunto finito de vectores es un cono convexo.
\end{proposition}

\begin{proof}
Tarea.
\end{proof}

\subsection{Factibilidad y conos generados por columnas}

Sea
\[
    A=
    \begin{pmatrix}
        \vert & \vert & & \vert\\
        a_1 & a_2 & \cdots & a_n\\
        \vert & \vert & & \vert
    \end{pmatrix}
    \in\mathbb{R}^{m\times n},
\]
donde \(a_1,\dots,a_n\in\mathbb{R}^m\) son las columnas de \(A\).

Observemos que
\[
    Ax
    =
    x_1a_1+x_2a_2+\cdots+x_na_n.
\]

Si \(x\geq 0\), entonces \(Ax\) es una combinación cónica de las columnas de \(A\). Por tanto,
\[
    C_A
    =
    \{Ax:x\geq 0\}
\]
es el cono generado por las columnas de \(A\).

Así, el sistema
\[
    Ax=b,\qquad x\geq 0
\]
es factible si y solamente si
\[
    b\in C_A.
\]

En consecuencia, el programa lineal en forma de igualdad
\[
    \begin{array}{ll}
        \max & c^\top x\\[1mm]
        \text{sujeto a} & Ax=b,\\
        & x\geq 0
    \end{array}
\]
es factible si y solamente si el vector \(b\) pertenece al cono generado por las columnas de \(A\).

% ============================================================
% APUNTES.pdf -- Trozo 3
% Páginas 5--6
% ============================================================

\subsection{Puntos extremos}

\begin{definition}[Punto extremo]
Sea \(X\subseteq \mathbb{R}^n\) un conjunto convexo. Un punto \(x_0\in X\) se llama \emph{punto extremo} de \(X\) si no puede escribirse como combinación convexa no trivial de dos puntos distintos de \(X\).

Es decir, si
\[
    x_0=tx_1+(1-t)x_2,
    \qquad t\in(0,1),
    \qquad x_1,x_2\in X,
\]
entonces necesariamente
\[
    x_1=x_2=x_0.
\]
\end{definition}

Geométricamente, los puntos extremos son los vértices del conjunto factible cuando este tiene estructura poliédrica.

\subsection{Soluciones básicas factibles}

Consideremos nuevamente un programa lineal en forma de igualdad:
\[
    \begin{array}{ll}
        \max & c^\top x\\[1mm]
        \text{sujeto a} & Ax=b,\\
        & x\geq 0,
    \end{array}
\]
donde
\[
    A\in\mathbb{R}^{m\times n},
    \qquad
    \operatorname{rango}(A)=m,
    \qquad
    b\in\mathbb{R}^m.
\]

Sea
\[
    [n]=\{1,\dots,n\}.
\]

Dado un subconjunto \(B\subseteq[n]\), denotamos por
\[
    A_B
\]
la matriz formada por las columnas de \(A\) cuyos índices pertenecen a \(B\).

Más explícitamente, si
\[
    A=
    \begin{pmatrix}
        \vert & \vert & & \vert\\
        a_1 & a_2 & \cdots & a_n\\
        \vert & \vert & & \vert
    \end{pmatrix},
\]
entonces
\[
    A_B
    =
    \begin{pmatrix}
        \vert & & \vert\\
        a_i & \cdots & a_j\\
        \vert & & \vert
    \end{pmatrix},
    \qquad i,j\in B.
\]

Análogamente, si \(x\in\mathbb{R}^n\), denotamos por
\[
    x_B
\]
el vector formado por las componentes \(x_i\) con índice \(i\in B\).

\begin{definition}[Base]
Sea \(A\in\mathbb{R}^{m\times n}\) con \(\operatorname{rango}(A)=m\). Un subconjunto
\[
    B\subseteq[n]
\]
se llama una \emph{base} si
\[
    |B|=m
\]
y la matriz
\[
    A_B
\]
es no singular.

Equivalente: \(B\) es una base si las columnas de \(A\) indexadas por \(B\) son linealmente independientes y forman una matriz cuadrada invertible.
\end{definition}

\begin{definition}[Solución básica]
Dado el programa lineal
\[
    \begin{array}{ll}
        \max & c^\top x\\[1mm]
        \text{sujeto a} & Ax=b,\\
        & x\geq 0,
    \end{array}
\]
sea \(B\subseteq[n]\) una base.

Un vector \(\widehat{x}\in\mathbb{R}^n\) se llama \emph{solución básica asociada a la base \(B\)} si satisface
\[
    A_B\widehat{x}_B=b
\]
y
\[
    \widehat{x}_j=0,
    \qquad
    \forall j\notin B.
\]

Las variables \(\widehat{x}_j\) con \(j\in B\) se llaman \emph{variables básicas}, mientras que las variables con \(j\notin B\) se llaman \emph{variables no básicas}.
\end{definition}

Como \(A_B\) es no singular, el sistema
\[
    A_B\widehat{x}_B=b
\]
tiene solución única, dada por
\[
    \widehat{x}_B=A_B^{-1}b.
\]

Por tanto, una vez escogida la base \(B\), queda determinada una única solución básica.

\begin{definition}[Solución básica factible]
Un vector \(\widehat{x}\in\mathbb{R}^n\) se llama \emph{solución básica factible} si es una solución básica asociada a alguna base \(B\subseteq[n]\) y, además, es factible para el programa lineal, es decir,
\[
    A\widehat{x}=b,
    \qquad
    \widehat{x}\geq 0.
\]
\end{definition}

En resumen:
\[
    \widehat{x}
    \text{ es SBF}
    \quad \Longleftrightarrow \quad
    \widehat{x}
    \text{ es solución básica y }
    \widehat{x}\geq 0.
\]

\subsection{Caracterización de soluciones básicas factibles}

Sea
\[
    x\in\mathbb{R}^n
\]
un punto factible, es decir,
\[
    Ax=b,
    \qquad
    x\geq 0.
\]

Definimos el conjunto de índices positivos de \(x\) por
\[
    K_x
    =
    \{j\in[n]:x_j>0\}.
\]

\begin{proposition}
Sea \(x\in\mathbb{R}^n\) un punto factible del sistema
\[
    Ax=b,
    \qquad
    x\geq 0.
\]
Entonces \(x\) es una solución básica factible asociada a alguna base \(B\subseteq[n]\) si y solamente si las columnas
\[
    \{a_j:j\in K_x\}
\]
son linealmente independientes.
\end{proposition}

\begin{proof}
Primero supongamos que \(x\) es una solución básica factible con base \(B\). Entonces, por definición,
\[
    x_j=0,
    \qquad
    \forall j\notin B.
\]
Por lo tanto, si \(x_j>0\), necesariamente \(j\in B\). Es decir,
\[
    K_x\subseteq B.
\]

Como \(B\) es una base, las columnas de \(A_B\) son linealmente independientes. Luego, cualquier subconjunto de ellas también es linealmente independiente. En particular, las columnas
\[
    \{a_j:j\in K_x\}
\]
son linealmente independientes.

Recíprocamente, supongamos ahora que las columnas
\[
    \{a_j:j\in K_x\}
\]
son linealmente independientes.

Si
\[
    |K_x|=m,
\]
tomamos
\[
    B=K_x.
\]
Entonces \(A_B\) es no singular y, como \(x_j=0\) para \(j\notin K_x\), el punto \(x\) es una solución básica factible con base \(B\).

Si
\[
    |K_x|<m,
\]
extendemos \(K_x\) a un conjunto \(B\subseteq[n]\) con
\[
    K_x\subseteq B,
    \qquad
    |B|=m,
\]
agregando índices de columnas de \(A\) de modo que las columnas correspondientes sigan siendo linealmente independientes. Esto es posible porque
\[
    \operatorname{rango}(A)=m.
\]

Entonces \(A_B\) es no singular. Además, como \(K_x\subseteq B\), se tiene
\[
    x_j=0,
    \qquad
    \forall j\notin B.
\]

Por otro lado, como \(x\) es factible,
\[
    b=Ax.
\]
Separando las columnas según \(B\) y su complemento, obtenemos
\[
    b=A_Bx_B+A_Nx_N,
\]
donde
\[
    N=[n]\setminus B.
\]
Pero \(x_N=0\), luego
\[
    b=A_Bx_B.
\]
Así, \(x\) es solución básica asociada a la base \(B\). Como además \(x\geq 0\), concluimos que \(x\) es solución básica factible.
\end{proof}

\begin{remark}
La proposición anterior dice que un punto factible es una solución básica factible precisamente cuando las columnas activas, es decir, aquellas asociadas a variables estrictamente positivas, son linealmente independientes.
\end{remark}

\subsection{Soluciones básicas factibles y puntos extremos}

\begin{proposition}
Toda solución básica factible es un punto extremo del conjunto factible
\[
    F=\{x\in\mathbb{R}^n:Ax=b,\ x\geq 0\}.
\]
\end{proposition}

\begin{proof}
Sea \(\widehat{x}\) una solución básica factible asociada a una base \(B\). Queremos probar que \(\widehat{x}\) no puede escribirse como combinación convexa no trivial de dos puntos factibles distintos.

Como \(\widehat{x}\) es solución básica asociada a \(B\), se tiene
\[
    \widehat{x}_j=0,
    \qquad
    \forall j\notin B,
\]
y
\[
    A_B\widehat{x}_B=b.
\]

Además, como \(A_B\) es no singular,
\[
    \widehat{x}_B=A_B^{-1}b.
\]

Supongamos que
\[
    \widehat{x}
    =
    t y+(1-t)z,
    \qquad
    t\in(0,1),
\]
con
\[
    y,z\in F.
\]

Como \(y,z\geq 0\) y \(\widehat{x}_j=0\) para todo \(j\notin B\), necesariamente
\[
    y_j=z_j=0,
    \qquad
    \forall j\notin B.
\]

Además, como \(y,z\in F\),
\[
    Ay=b,
    \qquad
    Az=b.
\]
Usando que sus componentes fuera de \(B\) son cero, se obtiene
\[
    A_By_B=b,
    \qquad
    A_Bz_B=b.
\]
Como \(A_B\) es invertible,
\[
    y_B=A_B^{-1}b=\widehat{x}_B,
    \qquad
    z_B=A_B^{-1}b=\widehat{x}_B.
\]
Por lo tanto,
\[
    y=z=\widehat{x}.
\]
Así, \(\widehat{x}\) es punto extremo de \(F\).
\end{proof}

\subsection{Número finito de soluciones básicas factibles}

Como una solución básica factible está asociada a una base \(B\subseteq[n]\) con \(|B|=m\), el número de posibles bases es finito. En particular,
\[
    \#\{\text{bases}\}
    \leq
    \binom{n}{m}.
\]

Por tanto,
\[
    \#\{\text{soluciones básicas factibles}\}
    \leq
    \binom{n}{m}
    <\infty.
\]

\begin{remark}
Esto es fundamental para el método simplex: aunque el conjunto factible puede tener infinitos puntos, los candidatos relevantes para optimizar una función lineal pueden buscarse entre un conjunto finito de soluciones básicas factibles.
\end{remark}

\subsection{Teorema fundamental de la programación lineal}

\begin{theorem}[Existencia de solución óptima básica factible]
Consideremos el programa lineal
\[
    \begin{array}{ll}
        \max & c^\top x\\[1mm]
        \text{sujeto a} & Ax=b,\\
        & x\geq 0.
    \end{array}
\]

Supongamos que el problema es factible y que la función objetivo está acotada superiormente sobre el conjunto factible. Es decir, existe \(M\in\mathbb{R}\) tal que
\[
    c^\top x\leq M,
    \qquad
    \forall x\in F.
\]

Entonces existe una solución óptima que es solución básica factible.
\end{theorem}

De manera equivalente, si el programa lineal tiene solución óptima, entonces posee al menos una solución óptima en un punto extremo del conjunto factible.

Como los puntos extremos corresponden a soluciones básicas factibles, el problema puede resolverse, al menos conceptualmente, revisando un número finito de candidatos:
\[
    \#\{\text{SBF}\}
    \leq
    \binom{n}{m}.
\]

\begin{remark}
Este resultado entrega la intuición central detrás del método simplex: buscar óptimos de una función lineal recorriendo soluciones básicas factibles, es decir, vértices del poliedro factible.
\end{remark}

\subsection{Flujos en redes}

Un ejemplo importante de programación lineal aparece en problemas de flujos sobre redes.

Sea
\[
    G=(V,E)
\]
un grafo dirigido, donde:
\begin{itemize}
    \item \(V\) es el conjunto de vértices;
    \item \(E\) es el conjunto de arcos o pares ordenados de vértices.
\end{itemize}

Un arco de \(E\) se escribe como
\[
    (a,s)\in E,
\]
donde el orden importa:
\[
    (a,s)\neq (s,a)
\]
en general.

\begin{remark}
En una red dirigida, cada arco tiene una orientación. Por eso los pares ordenados representan conexiones dirigidas entre vértices.
\end{remark}

% revisar: la página 6 introduce "Flujos en redes", pero el desarrollo
% completo del modelo parece continuar en la página siguiente.

% ============================================================
% APUNTES.pdf -- Trozo 4
% Páginas 7--8
% ============================================================

\begin{proof}[Demostración del teorema fundamental de la programación lineal]
Sea \(x^0\in F\) un punto factible. Consideremos el conjunto de puntos factibles cuyo valor objetivo es al menos tan bueno como el de \(x^0\):
\[
    F_{x^0}
    =
    \{x\in F:c^\top x\geq c^\top x^0\}.
\]

Como el problema está acotado superiormente, podemos escoger dentro de \(F_{x^0}\) un punto \(\widehat{x}\) que maximice el número de componentes nulas. Es decir, elegimos \(\widehat{x}\in F_{x^0}\) tal que
\[
    \#\{j:\widehat{x}_j=0\}
\]
sea máximo entre los puntos factibles con valor objetivo al menos \(c^\top x^0\).

Queremos probar que \(\widehat{x}\) es una solución básica factible.

Supongamos, por contradicción, que \(\widehat{x}\) no es una SBF. Definamos
\[
    K
    =
    \{j\in[n]:\widehat{x}_j>0\}.
\]
Como \(\widehat{x}\) no es solución básica factible, por la caracterización anterior las columnas
\[
    \{a_j:j\in K\}
\]
son linealmente dependientes.

Por lo tanto, existe un vector no nulo \(v_K\in\mathbb{R}^{K}\) tal que
\[
    A_Kv_K=0.
\]

Extendemos este vector a un vector \(v\in\mathbb{R}^n\) definiendo
\[
    v_j=0,
    \qquad
    j\notin K.
\]
Entonces
\[
    Av=0,
    \qquad
    v\neq 0.
\]

Además, como \(v_j=0\) para \(j\notin K\), las perturbaciones
\[
    \widehat{x}+tv
\]
solo modifican las coordenadas donde \(\widehat{x}\) es estrictamente positiva.

Para \(t\in\mathbb{R}\), se tiene
\[
    A(\widehat{x}+tv)
    =
    A\widehat{x}+tAv
    =
    b.
\]
Luego, mientras se mantenga la no negatividad,
\[
    \widehat{x}+tv\geq 0,
\]
el punto \(\widehat{x}+tv\) sigue siendo factible.

Ahora analizamos el valor objetivo:
\[
    c^\top(\widehat{x}+tv)
    =
    c^\top\widehat{x}
    +
    t\,c^\top v.
\]

\medskip

\noindent
\textbf{Caso 1: \(c^\top v>0\).}

Entonces, para \(t>0\) suficientemente pequeño,
\[
    \widehat{x}+tv\geq 0
\]
y
\[
    c^\top(\widehat{x}+tv)
    >
    c^\top\widehat{x}.
\]
Como además \(A(\widehat{x}+tv)=b\), esto produciría un punto factible con valor objetivo mayor.

Si el valor objetivo pudiera crecer indefinidamente en esa dirección, se contradiría la hipótesis de acotación superior. Por tanto, aumentando \(t\) hasta el mayor valor admisible, alguna coordenada positiva de \(\widehat{x}\) se anula y obtenemos un punto factible con:
\[
    c^\top(\widehat{x}+tv)\geq c^\top\widehat{x}
\]
y con más componentes nulas que \(\widehat{x}\), contradiciendo la elección de \(\widehat{x}\).

\medskip

\noindent
\textbf{Caso 2: \(c^\top v<0\).}

En este caso usamos la dirección contraria \(-v\). Como
\[
    c^\top(-v)>0,
\]
se aplica el caso anterior a \(-v\). Nuevamente obtenemos una contradicción con la maximalidad del número de ceros de \(\widehat{x}\).

\medskip

\noindent
\textbf{Caso 3: \(c^\top v=0\).}

Entonces
\[
    c^\top(\widehat{x}+tv)
    =
    c^\top\widehat{x}
\]
para todo \(t\) tal que \(\widehat{x}+tv\) sea factible.

Como \(v\neq 0\) y \(v\) solo tiene soporte en \(K\), podemos elegir \(t\neq 0\) de manera que:
\[
    \widehat{x}+tv\geq 0,
\]
y al menos una de las componentes positivas de \(\widehat{x}\) se vuelva cero.

Así obtenemos un punto factible \(\widetilde{x}\) tal que
\[
    c^\top \widetilde{x}
    =
    c^\top\widehat{x}
    \geq
    c^\top x^0,
\]
pero con más componentes nulas que \(\widehat{x}\), lo cual contradice nuevamente la elección de \(\widehat{x}\).

\medskip

En todos los casos llegamos a contradicción. Por lo tanto, las columnas asociadas a las componentes positivas de \(\widehat{x}\) deben ser linealmente independientes. Por la caracterización de soluciones básicas factibles, concluimos que
\[
    \widehat{x}
    \text{ es una SBF.}
\]

Como además \(\widehat{x}\) puede escogerse con valor objetivo óptimo, existe una solución óptima que es solución básica factible.
\end{proof}

\subsection{Interpretación geométrica del resultado}

Sea
\[
    F=\{x\in\mathbb{R}^n:Ax=b,\ x\geq 0\}\neq\varnothing.
\]

Si la función objetivo está acotada superiormente sobre \(F\), entonces existe una solución óptima que es solución básica factible.

La interpretación geométrica es:

\[
    \text{si un programa lineal tiene óptimo, entonces algún óptimo está en un vértice.}
\]

Como las soluciones básicas factibles corresponden a vértices del conjunto factible, el teorema justifica buscar el óptimo recorriendo SBF.

\subsection{Soluciones básicas factibles en forma general}

En la forma de igualdad
\[
    F=\{x\in\mathbb{R}^n:Ax=b,\ x\geq 0\},
\]
las restricciones de no negatividad
\[
    x_j\geq 0
\]
también pueden verse como restricciones activas o inactivas.

\begin{definition}[Restricción activa]
Para una solución factible \(x\), una restricción se dice \emph{activa} si se satisface con igualdad.
\end{definition}

En particular, para la restricción
\[
    x_j\geq 0,
\]
la restricción está activa en \(x\) si
\[
    x_j=0.
\]

\begin{definition}[Solución básica factible en términos de restricciones activas]
Una solución factible \(x\) de un programa lineal se llama \emph{solución básica factible} si entre las restricciones activas hay suficientes restricciones linealmente independientes para determinar de manera única al punto \(x\).
\end{definition}

Esta definición coincide con la definición anterior para la forma de igualdad.

\subsection{Ejemplos de restricciones activas}

\paragraph{Ejemplo 1.}
Consideremos un problema en una variable:
\[
    x\in\mathbb{R}.
\]
Si el conjunto factible está determinado por restricciones como
\[
    2x+5x_2=5,
\]
% revisar: en el manuscrito el ejemplo aparece abreviado; parece ser un ejemplo
% de restricción activa en dimensión baja.
entonces una restricción es activa cuando se alcanza exactamente la igualdad correspondiente.

\paragraph{Ejemplo 2.}
Consideremos restricciones del tipo
\[
    x_1+x_2=3.
\]
Un punto como
\[
    (1,0,1,0)
\]
puede representar una solución donde ciertas restricciones de no negatividad están activas, es decir, ciertas coordenadas son iguales a cero.

\begin{remark}
La idea general es que en una SBF hay suficientes restricciones activas para fijar completamente el punto. Geométricamente, esto corresponde a estar en un vértice del conjunto factible.
\end{remark}

\subsection{Forma de igualdad y soluciones básicas factibles}

Para un problema en forma de igualdad,
\[
    F=
    \{x\in\mathbb{R}^n:Ax=b,\ x\geq 0\},
\]
un punto \(x\in F\) es una SBF si existe una base
\[
    B\subseteq[n]
\]
tal que:
\[
    A_B \text{ es no singular},
\]
y
\[
    x_j=0,
    \qquad
    \forall j\notin B.
\]

Es decir, las variables no básicas son nulas.

\subsection{Caracterizaciones equivalentes de SBF}

\begin{theorem}
Sea
\[
    F=
    \{x\in\mathbb{R}^n:Ax=b,\ x\geq 0\}.
\]
Para un punto \(v\in F\), las siguientes condiciones son equivalentes:

\begin{enumerate}
    \item Existe un vector \(c\in\mathbb{R}^n\) tal que \(v\) es el único maximizador de \(c^\top x\) sobre \(F\), es decir,
    \[
        c^\top v>c^\top x,
        \qquad
        \forall x\in F,\ x\neq v.
    \]

    \item \(v\) es una solución básica factible.

    \item \(v\) es un punto extremo de \(F\), es decir, \(v\) no puede escribirse como combinación convexa no trivial de otros dos puntos de \(F\).
\end{enumerate}
\end{theorem}

\begin{proof}
Primero probamos que \((1)\Rightarrow(2)\).

Supongamos que existe \(c\in\mathbb{R}^n\) tal que
\[
    c^\top v>c^\top x,
    \qquad
    \forall x\in F,\ x\neq v.
\]
Entonces \(v\) es el único maximizador de \(c^\top x\) sobre \(F\).

Por el teorema fundamental de la programación lineal, existe una solución básica factible \(\widehat{x}\in F\) que maximiza \(c^\top x\) sobre \(F\). Luego,
\[
    c^\top \widehat{x}
    =
    c^\top v.
\]
Como \(v\) es el único maximizador, necesariamente
\[
    \widehat{x}=v.
\]
Por tanto, \(v\) es una SBF.

\medskip

Ahora probamos que \((2)\Rightarrow(1)\).

Supongamos que \(v\) es una solución básica factible asociada a una base \(B\subseteq[n]\). Definimos un vector \(c\in\mathbb{R}^n\) por
\[
    c_j=
    \begin{cases}
        0, & j\in B,\\
        -1, & j\notin B.
    \end{cases}
\]

Como \(v_j=0\) para todo \(j\notin B\), se tiene
\[
    c^\top v=0.
\]

Por otro lado, si \(x\in F\), entonces \(x\geq 0\), y por la definición de \(c\),
\[
    c^\top x
    =
    -\sum_{j\notin B}x_j
    \leq 0.
\]
Así, \(v\) maximiza \(c^\top x\) sobre \(F\).

Además,
\[
    c^\top x=0
\]
solo si
\[
    x_j=0,
    \qquad
    \forall j\notin B.
\]
En ese caso, como \(x\in F\),
\[
    Ax=b.
\]
Pero, al tener \(x_j=0\) para \(j\notin B\), esto se reduce a
\[
    A_Bx_B=b.
\]
Como \(A_B\) es no singular, la solución es única:
\[
    x_B=A_B^{-1}b=v_B.
\]
Por lo tanto,
\[
    x=v.
\]
Concluimos que
\[
    c^\top v>c^\top x,
    \qquad
    \forall x\in F,\ x\neq v.
\]

Luego \(v\) es el único maximizador de una función lineal sobre \(F\).

\medskip

La equivalencia entre \((2)\) y \((3)\) corresponde a la identificación geométrica entre soluciones básicas factibles y puntos extremos del poliedro factible. Ya probamos que toda SBF es punto extremo. Recíprocamente, si un punto extremo no fuera SBF, entonces existiría una dirección no trivial \(w\neq 0\) con
\[
    Aw=0
\]
y con soporte contenido en las coordenadas positivas de \(v\). Para \(t>0\) suficientemente pequeño,
\[
    v+tw\in F
    \qquad\text{y}\qquad
    v-tw\in F.
\]
Entonces
\[
    v
    =
    \frac{1}{2}(v+tw)
    +
    \frac{1}{2}(v-tw),
\]
lo que contradice que \(v\) sea punto extremo. Por tanto, todo punto extremo es una SBF.
\end{proof}

% ============================================================
% APUNTES.pdf -- Trozo 5
% Páginas 9--11
% ============================================================

\section{Linealización de funciones lineales a trozos convexas}

\subsection{Funciones lineales a trozos convexas}

Una clase importante de funciones que aparece en programación lineal es la clase de las funciones \emph{lineales a trozos convexas}.

\begin{definition}[Función lineal a trozos convexa]
Una función
\[
    f:\mathbb{R}^n\to\mathbb{R}
\]
se dice \emph{lineal a trozos convexa} si puede escribirse como el máximo de una cantidad finita de funciones afines, es decir, si existen vectores
\[
    c_1,\dots,c_m\in\mathbb{R}^n
\]
y escalares
\[
    d_1,\dots,d_m\in\mathbb{R}
\]
tales que
\[
    f(x)
    =
    \max_{i\in[m]}
    \{c_i^\top x+d_i\}.
\]
\end{definition}

Aquí usamos la notación
\[
    [m]=\{1,\dots,m\}.
\]

\begin{remark}
Toda función de la forma
\[
    f(x)=\max_{i\in[m]}\{c_i^\top x+d_i\}
\]
es convexa, porque es el máximo puntual de funciones afines.
\end{remark}

\subsection{Ejemplo: aproximación lineal}

Supongamos que tenemos puntos
\[
    p_i=(x_i,y_i)\in\mathbb{R}^2,
    \qquad i=1,\dots,m,
\]
y queremos encontrar una recta
\[
    y=ax+b
\]
que aproxime estos datos.

El error asociado al punto \(p_i\) es
\[
    r_i(a,b)
    =
    ax_i+b-y_i.
\]

Distintas medidas de aproximación conducen a distintos problemas de optimización.

\paragraph{Mínimos cuadrados.}
Una primera medida es minimizar la suma de los cuadrados de los residuos:
\[
    \min_{a,b}
    \sum_{i=1}^{m}
    (ax_i+b-y_i)^2.
\]

\paragraph{Error máximo.}
Otra posibilidad es minimizar el peor error absoluto:
\[
    \min_{a,b}
    \max_{i\in[m]}
    |ax_i+b-y_i|.
\]

\paragraph{Error absoluto total.}
También podemos minimizar la suma de los errores absolutos:
\[
    \min_{a,b}
    \sum_{i=1}^{m}
    |ax_i+b-y_i|.
\]

Estas dos últimas formulaciones involucran funciones lineales a trozos convexas y pueden transformarse en programas lineales mediante variables auxiliares.

\subsection{Ejemplos básicos}

\begin{example}[Valor absoluto]
La función
\[
    f:\mathbb{R}\to\mathbb{R},
    \qquad
    f(x)=|x|,
\]
puede escribirse como
\[
    |x|
    =
    \max\{x,-x\}.
\]
Por lo tanto, \(f\) es lineal a trozos convexa.
\end{example}

\begin{example}[Norma infinito]
La norma infinito en \(\mathbb{R}^n\) está dada por
\[
    \|x\|_\infty
    =
    \max_{i\in[n]}|x_i|.
\]
Como
\[
    |x_i|=\max\{x_i,-x_i\},
\]
se tiene
\[
    \|x\|_\infty
    =
    \max\{x_1,-x_1,\dots,x_n,-x_n\}.
\]
Luego \(\|\cdot\|_\infty\) es una función lineal a trozos convexa.
\end{example}

\begin{example}[Norma uno]
La norma uno en \(\mathbb{R}^n\) está dada por
\[
    \|x\|_1
    =
    \sum_{i=1}^{n}|x_i|.
\]
Como el valor absoluto es lineal a trozos convexo y las sumas con coeficientes no negativos preservan convexidad, \(\|\cdot\|_1\) es convexa.

Además, usando variables auxiliares, también puede representarse mediante restricciones lineales.
\end{example}

\subsection{Operaciones que preservan convexidad}

Sean
\[
    f_1,\dots,f_k:\mathbb{R}^n\to\mathbb{R}
\]
funciones convexas y sean
\[
    \alpha_1,\dots,\alpha_k\geq 0.
\]
Entonces
\[
    f(x)
    =
    \sum_{i=1}^{k}\alpha_i f_i(x)
\]
también es convexa.

En particular, si \(f_1,\dots,f_k\) son funciones lineales a trozos convexas, entonces cualquier combinación lineal no negativa de ellas sigue siendo convexa.

\subsection{Restricciones lineales a trozos convexas}

Supongamos que tenemos una restricción de la forma
\[
    \max_{i\in[m]}
    \{c_i^\top x+d_i\}
    \leq h.
\]

Esta restricción es equivalente a pedir que cada una de las funciones afines esté por debajo de \(h\). Es decir,
\[
    \max_{i\in[m]}
    \{c_i^\top x+d_i\}
    \leq h
\]
si y solamente si
\[
    c_i^\top x+d_i\leq h,
    \qquad i=1,\dots,m.
\]

Por tanto, una restricción lineal a trozos convexa puede reemplazarse por una familia finita de restricciones lineales.

\begin{remark}
Esta observación es fundamental: permite transformar restricciones con máximos de funciones afines en restricciones lineales ordinarias.
\end{remark}

\subsection{Linealización mediante epígrafo}

Sea
\[
    f:\mathbb{R}^n\to\mathbb{R}
\]
una función lineal a trozos convexa de la forma
\[
    f(x)
    =
    \max_{i\in[m]}
    \{c_i^\top x+d_i\}.
\]

Consideremos el problema
\[
    \begin{array}{ll}
        \min & f(x)\\[1mm]
        \text{sujeto a} & Ax\leq b.
    \end{array}
\]

La idea es introducir una nueva variable \(t\) que represente una cota superior para \(f(x)\). Así, el problema anterior es equivalente a
\[
    \begin{array}{ll}
        \min & t\\[1mm]
        \text{sujeto a} & t\geq f(x),\\
        & Ax\leq b.
    \end{array}
\]

Como
\[
    f(x)=\max_{i\in[m]}\{c_i^\top x+d_i\},
\]
la restricción
\[
    t\geq f(x)
\]
es equivalente a
\[
    t\geq c_i^\top x+d_i,
    \qquad i=1,\dots,m.
\]

Por tanto, obtenemos el programa lineal
\[
    \begin{array}{ll}
        \min & t\\[1mm]
        \text{sujeto a}
        & c_i^\top x+d_i\leq t,
        \qquad i=1,\dots,m,\\
        & Ax\leq b.
    \end{array}
\]

\begin{definition}[Epígrafo]
El \emph{epígrafo} de una función
\[
    f:\mathbb{R}^n\to\mathbb{R}
\]
se define como
\[
    \operatorname{epi}(f)
    =
    \{(x,t)\in\mathbb{R}^n\times\mathbb{R}:f(x)\leq t\}.
\]
\end{definition}

Con esta notación, la transformación anterior consiste en reemplazar la minimización de \(f(x)\) por la minimización de \(t\) sobre el epígrafo de \(f\).

\[
    \min_x f(x)
    \qquad
    \Longleftrightarrow
    \qquad
    \begin{array}{ll}
        \min & t\\
        \text{sujeto a} & (x,t)\in \operatorname{epi}(f).
    \end{array}
\]

\subsection{Linealización del valor absoluto}

La desigualdad
\[
    t\geq |z|
\]
es equivalente al sistema lineal
\[
    t\geq z,
    \qquad
    t\geq -z.
\]

En efecto,
\[
    |z|=\max\{z,-z\}.
\]

Por tanto, cada vez que aparece una expresión de la forma
\[
    |z|
\]
en una función objetivo o restricción, podemos introducir una variable auxiliar \(t\) y reemplazar
\[
    t\geq |z|
\]
por
\[
    t\geq z,
    \qquad
    t\geq -z.
\]

\subsection{Ejemplo: ajuste de una recta con norma infinito}

Sean
\[
    p_i=(u_i,v_i)\in\mathbb{R}^2,
    \qquad i=1,\dots,m.
\]

Queremos encontrar una recta
\[
    v=au+b
\]
que aproxime los puntos \(p_i\). El residuo asociado al punto \(p_i\) es
\[
    r_i(a,b)=au_i+b-v_i.
\]

Si medimos el error usando la norma infinito, el problema es
\[
    \min_{a,b}
    \max_{i\in[m]}
    |au_i+b-v_i|.
\]

Introducimos una variable auxiliar \(t\). El problema anterior es equivalente a
\[
    \begin{array}{ll}
        \min & t\\[1mm]
        \text{sujeto a}
        & t\geq |au_i+b-v_i|,
        \qquad i=1,\dots,m.
    \end{array}
\]

Usando la linealización del valor absoluto, esto se transforma en
\[
    \begin{array}{ll}
        \min & t\\[1mm]
        \text{sujeto a}
        & t\geq au_i+b-v_i,
        \qquad i=1,\dots,m,\\
        & t\geq -au_i-b+v_i,
        \qquad i=1,\dots,m.
    \end{array}
\]

Este es un programa lineal en las variables
\[
    a,\ b,\ t.
\]

\subsection{Ejemplo: ajuste de una recta con norma uno}

Ahora medimos el error total mediante la norma uno:
\[
    \min_{a,b}
    \sum_{i=1}^{m}
    |au_i+b-v_i|.
\]

Introducimos variables auxiliares
\[
    e_i\geq 0,
    \qquad i=1,\dots,m,
\]
tales que
\[
    e_i\geq |au_i+b-v_i|.
\]

La condición anterior es equivalente a
\[
    e_i\geq au_i+b-v_i,
    \qquad
    e_i\geq -au_i-b+v_i,
    \qquad i=1,\dots,m.
\]

Por tanto, el problema se escribe como el programa lineal
\[
    \begin{array}{ll}
        \min & \displaystyle\sum_{i=1}^{m}e_i\\[2mm]
        \text{sujeto a}
        & e_i\geq au_i+b-v_i,
        \qquad i=1,\dots,m,\\
        & e_i\geq -au_i-b+v_i,
        \qquad i=1,\dots,m,\\
        & e_i\geq 0,
        \qquad i=1,\dots,m.
    \end{array}
\]

\subsection{Ejemplo: planificación de producción de helados}

Consideremos un problema de planificación de producción durante un año, dividido en doce meses. Para cada mes \(i\in\{1,\dots,12\}\), la demanda está dada por
\[
    d_i.
\]

Definimos las variables:
\[
    X_i\geq 0,
    \qquad i=1,\dots,12,
\]
donde \(X_i\) representa la producción del mes \(i\), y
\[
    S_i\geq 0,
    \qquad i=1,\dots,12,
\]
donde \(S_i\) representa el stock al final del mes \(i\).

El balance de inventario se expresa como
\[
    X_i+S_{i-1}-S_i=d_i,
    \qquad i=1,\dots,12.
\]

Aquí \(S_{i-1}\) corresponde al inventario disponible al comienzo del mes \(i\).

En los apuntes se considera que el stock inicial y el stock final son cero:
\[
    S_0=0,
    \qquad
    S_{12}=0.
\]
% revisar: la frase manuscrita indica que el stock del año pasado y el stock final
% de este año deben ser cero.

La función objetivo incluye dos costos:
\begin{itemize}
    \item un costo de almacenamiento proporcional al stock:
    \[
        20000\sum_{i=1}^{12}S_i;
    \]
    \item un costo por cambios en la producción entre meses consecutivos:
    \[
        50000\sum_{i=1}^{12}|X_i-X_{i-1}|.
    \]
\end{itemize}

% revisar: el índice inicial de la suma de variaciones puede ser i=2,\dots,12
% si no se define una producción previa X_0. En el manuscrito aparece escrito
% como una variación entre meses consecutivos.

Así, una formulación inicial del problema es
\[
    \begin{array}{ll}
        \min
        &
        20000\displaystyle\sum_{i=1}^{12}S_i
        +
        50000\displaystyle\sum_{i=1}^{12}|X_i-X_{i-1}|
        \\[3mm]
        \text{sujeto a}
        &
        X_i+S_{i-1}-S_i=d_i,
        \qquad i=1,\dots,12,\\[1mm]
        &
        X_i\geq 0,
        \qquad i=1,\dots,12,\\[1mm]
        &
        S_i\geq 0,
        \qquad i=1,\dots,12,\\[1mm]
        &
        S_0=0,\qquad S_{12}=0.
    \end{array}
\]

Para linealizar los valores absolutos, introducimos variables auxiliares
\[
    t_i\geq 0,
    \qquad i=1,\dots,12,
\]
tales que
\[
    t_i\geq |X_i-X_{i-1}|.
\]

Equivalentemente,
\[
    t_i\geq X_i-X_{i-1},
    \qquad
    t_i\geq -X_i+X_{i-1}.
\]

Por lo tanto, obtenemos el programa lineal equivalente
\[
    \begin{array}{ll}
        \min
        &
        20000\displaystyle\sum_{i=1}^{12}S_i
        +
        50000\displaystyle\sum_{i=1}^{12}t_i
        \\[3mm]
        \text{sujeto a}
        &
        X_i+S_{i-1}-S_i=d_i,
        \qquad i=1,\dots,12,\\[1mm]
        &
        t_i\geq X_i-X_{i-1},
        \qquad i=1,\dots,12,\\[1mm]
        &
        t_i\geq -X_i+X_{i-1},
        \qquad i=1,\dots,12,\\[1mm]
        &
        X_i\geq 0,
        \qquad i=1,\dots,12,\\[1mm]
        &
        S_i\geq 0,
        \qquad i=1,\dots,12,\\[1mm]
        &
        t_i\geq 0,
        \qquad i=1,\dots,12,\\[1mm]
        &
        S_0=0,\qquad S_{12}=0.
    \end{array}
\]

\subsection{Ejemplo: separación de puntos}

Supongamos que tenemos dos familias de puntos en el plano:
\[
    p_1,\dots,p_m
\]
etiquetados con \(+1\), y
\[
    q_1,\dots,q_n
\]
etiquetados con \(-1\).

Escribimos
\[
    p_i=(u(p_i),v(p_i)),
    \qquad
    q_j=(u(q_j),v(q_j)).
\]

Queremos encontrar una recta
\[
    v=au+b
\]
que separe los puntos \(p_i\) de los puntos \(q_j\).

Una forma de imponer la separación es pedir que los puntos \(p_i\) queden por encima de la recta y los puntos \(q_j\) queden por debajo. Esto se expresa como
\[
    v(p_i)\geq au(p_i)+b,
    \qquad i=1,\dots,m,
\]
y
\[
    v(q_j)\leq au(q_j)+b,
    \qquad j=1,\dots,n.
\]

Estas son restricciones lineales en las variables \(a\) y \(b\).

\subsubsection{Separación con margen}

Podemos mejorar la formulación introduciendo una variable de margen
\[
    \delta\geq 0.
\]

La idea es exigir que los puntos \(p_i\) estén al menos a una distancia vertical \(\delta\) por encima de la recta, y que los puntos \(q_j\) estén al menos a una distancia vertical \(\delta\) por debajo de la recta.

Esto da las restricciones
\[
    v(p_i)\geq au(p_i)+b+\delta,
    \qquad i=1,\dots,m,
\]
y
\[
    v(q_j)\leq au(q_j)+b-\delta,
    \qquad j=1,\dots,n.
\]

Entonces buscamos maximizar el margen:
\[
    \begin{array}{ll}
        \max & \delta\\[1mm]
        \text{sujeto a}
        &
        v(p_i)\geq au(p_i)+b+\delta,
        \qquad i=1,\dots,m,\\[1mm]
        &
        v(q_j)\leq au(q_j)+b-\delta,
        \qquad j=1,\dots,n,\\[1mm]
        &
        \delta\geq 0.
    \end{array}
\]

Este es un programa lineal en las variables
\[
    a,\ b,\ \delta.
\]

\begin{remark}
Este ejemplo muestra cómo un problema geométrico de clasificación lineal puede escribirse como un programa lineal, siempre que usemos una noción lineal de separación y margen.
\end{remark}
